\subsection{Objetivo y contribuciones}

El objetivo principal de este trabajo es diseñar, desarrollar y evaluar un sistema 
descentralizado para el almacenamiento y verificación de ficheros, basado en la
integración de tecnologías de almacenamiento distribuido y registro inmutable.
En particular, se propone una solución que permite almacenar documentos en 
IPFS y registrar sus identificadores criptográficos en una red blockchain, 
con el fin de garantizar su integridad y verificabilidad.

A diferencia de soluciones que emplean redes públicas, en este trabajo 
se utiliza una infraestructura blockchain local basada en Hardhat, 
lo que permite un entorno de experimentación controlado y reproducible. 
Este enfoque facilita el análisis del comportamiento del sistema sin las 
limitaciones derivadas de costes económicos o latencias variables propias 
de redes públicas, centrando el estudio en aspectos técnicos como el rendimiento 
y la eficiencia.

El sistema desarrollado incluye una interfaz de usuario basada en tecnologías
web (HTML, CSS y JavaScript) que permite interactuar con ambas capas del 
sistema: almacenamiento en IPFS y registro en blockchain. A través de esta interfaz,
el usuario puede almacenar ficheros, recuperar información 
y verificar la autenticidad de los datos mediante la comparación de hashes.

Además del desarrollo del prototipo, este trabajo define un conjunto de 
métricas orientadas a evaluar el comportamiento del sistema. 
Entre ellas se incluyen el tiempo de almacenamiento, el tiempo de recuperación, 
el impacto del tamaño del fichero en el rendimiento y el coste temporal añadido
por el uso de blockchain. Estas métricas permiten realizar una comparación 
objetiva entre el uso exclusivo de IPFS y el enfoque híbrido propuesto.

Las principales contribuciones de este trabajo pueden resumirse en los 
siguientes puntos. (I) Diseño e implementación de un sistema
funcional que integra almacenamiento descentralizado y verificación basada 
en blockchain. (II) Definición de un marco experimental que
permite medir y comparar el rendimiento de ambas tecnologías. (III) Análisis crítico de los resultados obtenidos, identificando las ventajas,
limitaciones y compromisos asociados al uso combinado de IPFS y 
blockchain en escenarios reales.

En conjunto, este trabajo proporciona una base experimental que 
permite analizar de forma sistemática las ventajas y limitaciones 
de los enfoques descentralizados actuales, contribuyendo a una 
mejor comprensión de su aplicabilidad en sistemas reales.